\begin{tikztimingtable}[
    timing/slope=0.05,
    timing/coldist=4pt,
    xscale=2.2,
    yscale=1.1,
    timing/dslope=0.1
]
  \texttt{clk}            & 16{C} \\
  \texttt{datain[3:0]}    & 2D{5} 2D{3} 12D{0} \\
  \texttt{constin[3:0]}   & 8D{0} 2D{2} 6D{0} \\
  \texttt{control[13:0]}  & 2D{00A0} 2D{10A0} 2D{2492} 2D{3B00} 2D{2ED5} 2D{3B00} 2D{2A81} 2D{3B00} \\
  \texttt{dataout[3:0]}   & 1X 3D{5} 2D{3} 2D{8} 2D{2} 3D{6} 3D{12} \\
  \texttt{addressout[3:0]}& 1X 5D{5} 2X 1D{8} 1D{6} 2X 1D{6} 1D{12} 2X \\
  \texttt{flags[3:0]}     & 1X 3D{0000} 2D{0100} 2X 2D{0001} 2X 1D{0000} 1D{0100} 2X \\
\extracode
  % Fondo sombreado alternado para distinguir cada ciclo de reloj
  \begin{pgfonlayer}{background}
    \foreach \t in {0, 4, 8, 12} {
      \pgfmathsetmacro{\tend}{\t+2}
      \fill[blue!4] (\t, 2.2) rectangle (\tend, -14.5);
    }
    \vertlines[gray!40, dashed, line width=0.5pt]{0, 2, ..., 16}
  \end{pgfonlayer}

  % Tiempos numéricos (eje temporal superior, cada 20 ns)
  \foreach \i in {0, ..., 8} {
    \pgfmathsetmacro{\x}{\i * 2}
    \pgfmathsetmacro{\tval}{\i * 20}
    \node[font=\sffamily\tiny, text=gray!80!black, anchor=south] at (\x, 4.4) {$\pgfmathprintnumber{\tval}$\,ns};
  }

  % Ciclos T0 a T7 y descripción de operaciones
  \foreach \i/\desc in {
    0/{Carga $R_0$},
    1/{Carga $R_1$},
    2/{Suma},
    3/{Lectura},
    4/{Resta},
    5/{Lectura},
    6/{SHL},
    7/{Lectura}
  } {
    \pgfmathsetmacro{\xc}{\i * 2 + 1}
    \node[font=\sffamily\scriptsize\bfseries, blue!80!black, anchor=south] at (\xc, 3.4) {$T_\i$};
    \node[font=\sffamily\tiny, text=darkgray, anchor=south] at (\xc, 2.4) {\desc};
  }
\end{tikztimingtable}
